\documentclass[11pt,a4paper]{article}

\usepackage[margin=1in]{geometry}
\usepackage{amsmath,amssymb,amsthm}
\usepackage{booktabs}
\usepackage{xcolor}
\usepackage[colorlinks=true,linkcolor=blue!60!black,citecolor=blue!60!black,
            urlcolor=blue!60!black]{hyperref}
\usepackage{microtype}
\usepackage{enumitem}
\usepackage[most]{tcolorbox}

\newtheorem{theorem}{Theorem}
\newtheorem{proposition}[theorem]{Proposition}
\theoremstyle{definition}
\newtheorem{definition}[theorem]{Definition}

\newcommand{\conj}[1]{\overline{#1}}
\newcommand{\lean}[1]{\texttt{\small #1}}
\newtcolorbox{keybox}[1][]{colback=blue!3,colframe=blue!45!black,
  boxrule=0.5pt,left=4pt,right=4pt,top=3pt,bottom=3pt,#1}
\newtcolorbox{retraction}[1][]{colback=red!3,colframe=red!45!black,
  boxrule=0.5pt,left=4pt,right=4pt,top=3pt,bottom=3pt,#1}

\title{\textbf{Formalizing quantum-fluid structure in Lean~4:}\\
\large a machine-checked library, what the formalization caught,\\ and what it could not}

\author{Xavier Callens\\
\small SocrateAI-Scientific-QuantumFluids\\
\small \url{https://github.com/xaviercallens/SocrateAI-Scientific-QuantumFluids}}
\date{September 2026}

\begin{document}
\maketitle

\begin{abstract}
\noindent
We present an open Lean~4/Mathlib library of 113 machine-checked theorems on the mathematical
structure of quantum fluids: correctness of the phase-winding algorithm by which every
Gross--Pitaevskii code locates vortices, quantization of circulation in the discrete form in
which it is actually computed, sign and conservation structure of the Galerkin-truncated
Gross--Pitaevskii nonlinearity on an arbitrary mode set, the Madelung decomposition, and a
bridge that restates these objects inside the vocabulary of OpenAI's Lean formalization of
the Navier--Stokes problem, imported as a real dependency. Every theorem carries a generated
axiom audit and is independently re-checked against two proof kernels.

\textbf{This is not a physics-discovery paper, and we say why.} The project began as one. Two
results we had regarded as findings were withdrawn after a literature check: a cubic
invariant of a complexified dyadic shell model is a special case of the Hamiltonian
structure of Vladimirova, Shavit and Falkovich and of L'vov, Podivilov and Procaccia; and a
``dual length'' bound on the Bogoliubov dispersion is dimensionally forced and weaker than
Onsager's inequality where it is not vacuous. We report these retractions in full, because
they are the clearest data we have on the actual subject of the paper: \emph{what formal
verification does and does not buy in physics}. It caught a false lemma we were about to
state, an inequality our prose had turned into an equality, a constant that made a bound
useless in the limit that mattered, eleven theorems that had silently escaped audit, and a
lemma whose \texttt{private} modifier had exempted it from independent checking. It
could not tell us that a correct theorem was already known, nor that a hypothesis was
physically naive; a literature search and a simulation did that. We give the library to the
community as verification infrastructure, with the open formalization targets we consider
most valuable.
\end{abstract}

\section{Introduction}

Proof assistants are now used at scale on fluid equations: OpenAI's announced constructions
for the Navier--Stokes and Euler equations ship with a Lean~4 development of several
thousand files~\cite{OpenAINSE}. Quantum fluids---superfluid helium, atomic Bose--Einstein
condensates---have, to our knowledge, no comparable formal footprint, although they are in
one respect better suited to it: circulation is quantized, the healing length is fixed by
constants, and the central algorithms of the field (vortex detection by phase winding,
Fourier--Galerkin truncation) are exactly discrete, so their correctness statements are
finite and provable rather than asymptotic.

This paper contributes a library, and a candid account of building it. The account matters
as much as the library. Over roughly six weeks the project produced, alongside correct
proofs, a sequence of claims that did not survive: two that a literature check identified as
rediscoveries or dimensionally forced (\S\ref{sec:limits}), and a topological hypothesis
that a well-resolved simulation contradicts (\S\ref{sec:empirical}). Each failure was caught
by a mechanism that was in place \emph{before} the claim was tested---a pre-registered
control, an external review, a literature gate. We think the distribution of which
mechanism caught which error is informative for anyone using proof assistants in the
physical sciences, and \S\ref{sec:caught}--\ref{sec:limits} are organised around it.

\begin{keybox}
\textbf{Stance.} No novelty is claimed for any physical statement in this paper. The
contribution is (i) the machine-checked library and its verification tooling, (ii) one
correctness result for a standard algorithm that we have not found stated elsewhere with its
failure case (\S\ref{sec:winding}), and (iii) the error taxonomy.
\end{keybox}

\section{The library}
\label{sec:library}

\paragraph{Build.} Lean~4.34.0-rc2 with Mathlib at tag \texttt{v4.34.0-rc2}. A single
\lean{import QuantumFluids} brings in fourteen modules (Table~\ref{tab:modules}); the full
build is 8792 jobs, zero errors, zero \lean{sorry} in any default target.

\begin{table}[h]
\centering\small
\begin{tabular}{@{}llp{8.6cm}@{}}
\toprule
module & thms & content \\
\midrule
\lean{VortexWinding} & 9 & loop quantization of the principal-branch phase sum; edge
  cancellation and its exact failure case; the oriented sum over a closed surface vanishes, so
  vortex lines do not end \\
\lean{QuantizedCirculation} & 6 & $\Gamma = q\,\kappa$, $\kappa = h/m$; no fraction of a
  quantum; the quantum is attained \\
\lean{GPGalerkin} & 13 & truncated Gross--Pitaevskii on any finite mode set: positivity of the
  interaction sum, mass and energy algebra, Hamiltonian gradient identity \\
\lean{HeliumKinematics} & 15 & identities a neutron-scattering analysis of $^4$He rests on:
  three-phonon decay open \emph{iff} the dispersion is anomalous, two-roton momentum range,
  calibration-invariance of a comparison with $2\Delta_R$, time-of-flight equivalences \\
\lean{BoseIntegral} & 7 & $\int_0^\infty t^{n}/(e^t-1)\,dt = n!\,\zeta(n+1)$ for every $n \ge 1$,
  with the even case explicit via Bernoulli numbers; the Debye $T^3$ coefficient \\
\lean{MadelungSplit} & 6 & $|\nabla\psi|^2 = |\nabla a|^2 + a^2|\nabla S|^2$ \\
\lean{MadelungNSE} & 2 & the Madelung velocity in the vocabulary of~\cite{OpenAINSE};
  $\nabla\!\cdot\!\nabla\varphi = \Delta\varphi$ in their definitions \\
\lean{DualLength} & 13 & algebra of $\varepsilon^2/(\hbar^2c^2k^3)$ for the Bogoliubov
  dispersion (significance withdrawn, \S\ref{sec:limits}) \\
\lean{QuantumFluidsShell} & 12 & complexified dyadic shell model: conservation, invariant
  real subspace, trace identities \\
\lean{ShellHamiltonian}, \lean{SigmaRule} & 15 & its cubic invariant and cutoff-uniform
  bounds (novelty withdrawn, \S\ref{sec:limits}) \\
\lean{Duality}, \lean{RipsFloor} & 15 & auxiliary inequalities; the 1-skeleton lemma used by
  the topological analysis \\
\bottomrule
\end{tabular}
\caption{Modules. Axiom footprint of all 113 theorems:
$\{\lean{propext}, \lean{Classical.choice}, \lean{Quot.sound}\}$.}
\label{tab:modules}
\end{table}

\paragraph{Verification stack.} Three layers, each added after a specific failure.
\begin{enumerate}[leftmargin=*,itemsep=2pt]
\item \emph{Generated axiom audit.} A script regenerates the \lean{\#print axioms} block of
  every file from its declarations and has a \lean{--check} mode. It replaced hand-kept
  lists after an external review found those had drifted (\S\ref{sec:caught}, item~5).
\item \emph{Comparator}~\cite{Comparator}. For each module a challenge file---definitions
  verbatim, proofs replaced by \lean{sorry}---is generated from the solution, and Comparator
  checks statement identity, an axiom whitelist, and acceptance by both the Lean kernel and
  the independent \lean{nanoda} kernel, inside a sandbox. Because challenges are generated
  from solutions, statement identity holds by construction; the independent content is the
  second kernel and the whitelist. Whether a statement is the \emph{right} statement remains
  a human audit, and \S\ref{sec:caught} shows why that caveat is not a formality.
\item \emph{Negative controls.} For each headline theorem a perturbed statement (a wrong
  constant, a dropped conjugate, a flipped inequality) is compiled and required to fail. A
  checker that has never been seen to reject is not yet a checker.
\end{enumerate}

\section{Selected results}

\subsection{Correctness of phase-winding vortex detection}
\label{sec:winding}

Every Gross--Pitaevskii simulation code locates quantized vortices the same way: sum the
principal-branch phase differences of $\psi$ around a grid plaquette, divide by $2\pi$,
round. Let $\mathrm{pd}(a,b)\in(-\pi,\pi]$ denote the principal representative of $b-a$.

\begin{theorem}[\lean{loop\_sum\_eq\_mul}]
For any phases $\theta_0,\dots,\theta_n$ with $\theta_n=\theta_0$,
$\sum_{i<n}\mathrm{pd}(\theta_i,\theta_{i+1}) \in 2\pi\mathbb{Z}$.
\end{theorem}
So the computed winding number is an integer by theorem, not by numerical accident. The
second fact codes rely on---that an edge shared by two plaquettes contributes oppositely to
each, so windings add like a discrete Stokes theorem---is usually stated as antisymmetry,
$\mathrm{pd}(a,b) = -\mathrm{pd}(b,a)$. \textbf{That statement is false}, and attempting to
prove it is how we found out:

\begin{theorem}[\lean{pdiff\_add\_rev}, \lean{pdiff\_add\_rev\_eq\_two\_pi\_iff}]
$\mathrm{pd}(a,b)+\mathrm{pd}(b,a)\in\{0,\,2\pi\}$, and the value $2\pi$ occurs exactly when
$\mathrm{pd}(a,b)=\mathrm{pd}(b,a)=\pi$.
\end{theorem}
At a phase step of exactly $\pi$ both traversal directions return $+\pi$, cancellation
between neighbouring plaquettes fails, and a spurious unit of winding can appear. The
folklore criterion ``resolve the field so that no phase step reaches $\pi$'' is thereby a
theorem with a sharp boundary. We have not found the disjunctive form stated in the
literature; it is elementary, and we claim only that it is now checked.

\subsection{Quantized circulation}

With superfluid velocity $u=(\hbar/m)\nabla S$, define the discrete circulation around a
closed loop of sample points as $\Gamma = (\hbar/m)\sum_i \mathrm{pd}(S_i,S_{i+1})$ and
$\kappa = 2\pi\hbar/m$.

\begin{theorem}[\lean{circulation\_quantized}, \lean{abs\_circulation\_ge},
\lean{circulation\_quantum\_attained}]
$\Gamma = q\kappa$ for an integer $q$; if $\Gamma\neq0$ then $|\Gamma|\ge\kappa$; and an
explicit four-point loop attains $\Gamma=\kappa$.
\end{theorem}
The attained case is included deliberately: a lower bound with no witness can be vacuous,
and the witness's closing step wraps by $2\pi$, exercising exactly the branch behaviour the
definition exists to handle. \emph{Not} proved: that the discrete sum equals the continuum
line integral. That is a sampling statement, and by the previous subsection it fails, by
design, when a phase step reaches $\pi$.

\subsection{Truncated Gross--Pitaevskii on an arbitrary mode set}

Let $G$ be an additive commutative group ($\mathbb{Z}^3$ in applications), $\Lambda\subset G$
finite, and $N_k=\sum_{k_1+k_3=k+k_2}\psi_{k_1}\conj{\psi_{k_2}}\psi_{k_3}$ the projected cubic
nonlinearity.

\begin{theorem}[\lean{pairing\_eq\_sum}, \lean{pairing\_nonneg}, \lean{grad\_identity},
\lean{energy\_rate\_zero}]
For every finite $\Lambda$: $Q:=\sum_{k\in\Lambda}\conj{\psi_k}N_k=\sum_q|A_q|^2\ge0$ with
$A_q=\sum_{a-b=q}\psi_a\conj{\psi_b}$; $\;\sum_k\mathrm{Im}(\conj{\psi_k}G_k)=0$ for
$G_k=\omega_k\psi_k+gN_k$; the polarised gradient identity
$\sum_q 2\mathrm{Re}(\conj{A_q}\,dA_q[\delta]) = 4\,\mathrm{Re}\sum_k\conj{\delta_k}N_k$
holds; and the first variation of $E=\sum\omega_k|\psi_k|^2+\tfrac g2Q$ along $\delta=-iG$
vanishes.
\end{theorem}
The statements are uniform in the truncation: no property of $\Lambda$ beyond finiteness is
used, which is the form in which such identities are needed when a cutoff is to be removed.
What is proved is the algebraic core; the chain-rule step identifying these variations with
time derivatives along a solution is not formalised, and the module header says so.

\subsection{A cross-project import}

\lean{MadelungNSE} imports \lean{NavierStokes.ProblemStatement} from~\cite{OpenAINSE} as a
Lake dependency pinned to a commit whose headline theorems had been independently audited.
In \emph{their} definitions of divergence and pressure gradient we prove
$\nabla\!\cdot\!\nabla\varphi=\Delta\varphi$---absent from their tree, where pressure enters
only through its gradient---and hence $\nabla\!\cdot u=(\hbar/m)\Delta S$ for the Madelung
velocity. The mathematics is elementary. The point is operational: the import became
possible only when both projects were aligned on one toolchain and one Mathlib revision;
before that, version skew, not missing theorems, was the entire obstacle to reuse. One
convenience lemma was dropped rather than carried with a \lean{sorry} when it did not
survive a coercion change in Mathlib; nothing depended on it.

\section{What the formalization caught}
\label{sec:caught}

Each item is an error that existed in prose, code, or intent, and was surfaced by the act of
formalizing or by the verification stack.

\begin{enumerate}[leftmargin=*,itemsep=3pt]
\item \textbf{A false lemma, before it was used.} Antisymmetry of the principal phase
  difference (\S\ref{sec:winding}). Our own extractor's documentation asserted it.
\item \textbf{An inequality that prose had made an equality.} A design memo described a
  structure-factor bound as holding ``exactly'' under Feynman's relation. The Lean theorem
  could only be stated with the Bijl--Feynman relation as an \emph{inequality} hypothesis,
  $\varepsilon\le\hbar^2k^2/(2mS)$, which is what it is. The formal statement was right
  while the accompanying prose was wrong; the discrepancy was found when the two were
  compared.
\item \textbf{A constant that voids a bound where it matters.} A cutoff-uniform enstrophy
  bound for a third-order dispersive regulator, $2D\Omega\le H+(2E)^{3/2}$, looked relevant
  to a uniformity question in a sibling project. Written out for the formal statement, the
  bound on $\Omega$ carries $1/D$ and diverges as the regulator is removed---precisely the
  class of bound that project had already ruled out. We had described the result as
  possibly significant; the fully explicit statement is what retracted that.
\item \textbf{Hypotheses that locate the physics.} The Madelung identities require
  differentiability of amplitude and phase at the point. That hypothesis is not
  bureaucratic: it fails exactly on a vortex line, which is where the hydrodynamic picture
  itself fails. The formal statement makes the domain of validity part of the theorem.
\item \textbf{Audit drift.} An external reviewer rebuilding the tagged release reported 13
  theorems in a file documented as having 12. Checking every file showed 76 theorems against
  65 audit lines: eleven theorems had never been axiom-audited, and every count we had
  published conflated ``theorems'' with ``theorems carrying an audit line''. All eleven were
  clean; the defect was in the reporting. The generated audit (\S\ref{sec:library}) is the fix.
\item \textbf{A lemma silently exempt from independent checking.} One helper was declared
  \lean{private}. A private name is mangled, so it cannot be exported to the second kernel:
  the run aborted rather than passing, and only then did we notice that the lemma had never
  been independently checked. It is now public. A visibility modifier had been acting as an
  unintended exemption from verification.
\item \textbf{A clone-breaking artefact.} The same review found the build directory
  committed as a symlink to a local disk, because an ignore pattern with a trailing slash
  matches directories and not symlinks. No proof assistant catches this; a second person
  building from the tag did.
\end{enumerate}

\section{What it could not catch}
\label{sec:limits}

\begin{retraction}
\textbf{Two results were correct, machine-checked, and not contributions.} The kernel
certifies that a statement follows from its hypotheses. It is silent on whether the
statement is new, and on whether it could have come out otherwise.
\end{retraction}

\paragraph{A rediscovery.} We had found that the complexified dyadic shell model
$\dot v_n = k_{n-1}v_{n-1}^2-k_n\conj{v_n}v_{n+1}$ becomes Hamiltonian under
$w_n=2^{-n/2}v_n$, with a cubic invariant
$H=\sum_n2^{-n}[\omega_n|v_n|^2+k_n\,\mathrm{Im}(\conj{v_n}^2v_{n+1})]$, energy conservation
as a Manley--Rowe relation, and volume preservation as a corollary; we proved it in Lean
and confirmed by a nullspace search that the invariants are exactly
$\{\text{mass},\text{mass}^2,H\}$. All of this is a degenerate case of published
structure: the canonical Hamiltonian form of the Sabra model under the analogous rescaling
and its cubic invariant~\cite{Lvov1999,Ditlevsen2000}; the three-wave chain with graded
phase symmetry and Manley--Rowe families~\cite{Vladimirova2021PRX}; second-harmonic
generation as a minimal turbulence model~\cite{Vladimirova2021PRE}, whose title is the
framing we had taken for an insight; and the two-mode invariant itself~\cite{ABDP1962}.
Phase-space volume preservation is a stated \emph{design criterion} for shell
models~\cite{Biferale2003}. The dyadic case is, if anything, poorer: $2^a+2^b=2^c$ forces
$a=b$, so its Manley--Rowe family has one member where the Fibonacci lattice has two.

\paragraph{A dimensionally forced bound.} For the Bogoliubov dispersion, which is the sum of
squares of its phonon and free-particle branches~\cite{Dalfovo1999},
$\ell(k)=\varepsilon^2/(\hbar^2c^2k^3)$ equals $1/k+k/k_*^2$ with $k_*=2mc/\hbar$, is
invariant under $k\mapsto k_*^2/k$, and is bounded below by $\sqrt2\,\xi$. We proved this,
and measured that superfluid $^4$He violates the bound by factors of 21 to 51 across seven
pressures using the published dispersion of~\cite{Godfrin2021}. But a theory with one length
and one velocity admits no other outcome: any such length with an interior minimum has it
near $1/\xi$ with value near $\xi$. The implied structure-factor bound
$S(k)\le\sqrt{k/2k_*}$ is weaker than Onsager's inequality $S(k)\le\hbar k/2mc$
\cite{Wreszinski2005,Stringari1993} for $k<k_*/2$ and vacuous beyond $2k_*$. And since
$\ell=v_L^2/(c^2k)$ with $v_L=\varepsilon/\hbar k$, any fluid with a roton violates the
bound: the measurement confirms that $^4$He has a roton. A proposed test on a dilute
condensate~\cite{Steinhauer2002} was therefore not run; its outcome is forced.

\paragraph{The lesson we draw.} Formal verification moved our error rate on
\emph{internal} questions---is this statement true, are its hypotheses complete, is this
constant where we think it is---close to zero. It did nothing for the \emph{external}
questions: is it known, is it forced, does it matter. Both retractions above were obtained
by a literature search we had listed as a blocking gate and then ran late. In a field where
the mathematics is a century deep, that gate should come first; a proof assistant makes it
cheaper to be rigorously unoriginal.

\section{Empirical companion work}
\label{sec:empirical}

Alongside the library we pursued a topological test: quantized circulation with a finite
core suggests a minimum vortex separation near $\xi$, which would appear as a floor in the
$H_0$ persistence of the vortex configuration (computed with GUDHI~\cite{GUDHI}; the only
formal ingredient is the 1-skeleton lemma of \lean{RipsFloor}). The sequence of outcomes is
instructive mainly about controls.

\begin{enumerate}[leftmargin=*,itemsep=3pt]
\item On a published $256^3$ Gross--Pitaevskii snapshot~\cite{PolancoData}, proximity-based
  line segmentation gave a floor of $1.49\,\xi$ against a null of $0.10\,\xi$. A sweep of the
  segmentation threshold showed the floor tracking the threshold (ratio $1.00$--$1.24$ over
  a fourfold change): the statistic measured its own parameter. Withdrawn.
\item Threshold-free tracing by cube adjacency (no tunable parameter; zero ambiguous cubes
  in $46\,762$) gave $0.94\,\xi$, six times the null---but a second pre-registered criterion,
  the fraction of separations below $\xi$, was indistinguishable from the null. Reported as a
  split verdict; that dataset resolves $\xi$ by only $1.5$ grid cells.
\item A survey found no public dilute-condensate dataset resolving $\xi$ by five cells or
  more (best: $2.3$), so we wrote a split-step solver, validated by energy convergence, core
  formation at the correct size, and dipole annihilation, and ran at $\xi/\Delta x=8$.
  Two properties of the method had to be learned and are now pinned as tests:
  imaginary-time relaxation first forms vortex cores and then \emph{destroys} the vortices,
  and the timestep must respect $\Delta t\,k_{\max}^2/2\ll1$ (a run fifty times over the
  bound is detectably wrong). At $t=20$ the
  minimum separation is $0.28\,\xi$ ($\Delta t=0.01$) and $0.38\,\xi$ ($\Delta t=0.005$),
  inside the matched Poisson band ($0.09$--$0.78\,\xi$) and far below $\xi$; and the fraction
  of separations below $\xi$, $0.07$--$0.08$, lies \emph{above} the null's 95th percentile
  ($0.04$). Vortices sit \emph{closer} than random placement, not farther. Bulk statistics are
  converged in the timestep (count $3\%$, mean separation $2\%$); the minimum, an extreme
  statistic of a chaotic system, is not ($34\%$), and the energy drift ($2.4\%$) still
  exceeds the solver's own strict criterion, so the refutation rests on the verdict being
  identical at both timesteps rather than on any single number. \textbf{By the
  pre-registered criteria the hypothesis is refuted.}
\end{enumerate}
The hypothesis was physically naive. Opposite-sign vortices attract and annihilate, a process
that passes continuously through separations below $\xi$; a snapshot of a live configuration
therefore contains pairs in mid-approach. Quantization constrains circulation, not distance.

\section{Applying the library to published experimental work}
\label{sec:godfrin}

The clearest use we have found for a library of this kind is not to generate physics but to check
derivations that an experimental paper states without proof. We took as subject the two publications
of Godfrin and co-workers in 2021--2026: the pressure-dependent dispersion of superfluid $^4$He from
inelastic neutron scattering~\cite{Godfrin2021} and its companion review.

\lean{HeliumKinematics} formalises relations that paper's analysis leans on: three-phonon decay is
kinematically open \emph{iff} the dispersion is anomalous, with excess energy exactly
$3ca\,k_1k_2(k_1+k_2)$; two rotons carry total momentum in $[0,2k_R]$, attained by parallel and
anti-parallel pairs, which is the review's statement that the Pitaevskii plateau "ends at $2k_\Delta$"
and "can in principle extend to zero wave number"; every time-of-flight energy is proportional to the
incident energy, so \emph{any} comparison of $\varepsilon(k)$ with $2\Delta_R$ is invariant under
recalibration of the roton gap; and Eq.~(7) of~\cite{Godfrin2021} is $dP/d\rho$ of its Eq.~(6).

Two consequences are worth reporting. First, with the published low-$k$ coefficients each closing
condition factors as $ck^n$ times a \emph{quadratic}, so the thresholds are explicit: soft emission
closes at $0.404\,\text{\AA}^{-1}$, the symmetric split at $0.455$, the phase velocity reaches $c$ at
$0.566$. These are three different numbers, often quoted as one $k_c$, and the symmetric split is the
last two-phonon channel to close; the measured table alone, with no fitted series, gives $0.453$.
Second, \lean{BoseIntegral} verifies the leading coefficient of that paper's specific-heat series ---
the series for which it notes that earlier published versions "contain errors" --- from
$\int_0^\infty t^3/(e^t-1)\,dt = \pi^4/15$ through to $A = 2\pi^2k_B^4V/(15c^3\hbar^3) =
0.0831\,\mathrm{J/(mol\,K^4)}$, matching the printed value. Mathlib had $\zeta(4)$ and the Gamma
integral but not the Bose integral.

We also record two apparent misprints in the preprint, demonstrated numerically against the paper's
own tables rather than asserted, in \texttt{docs/FOR\_GODFRIN.md}. That note has not been sent; it is
a draft for the repository owner.

\section{A proposal the literature gate stopped}
\label{sec:gate}

\S\ref{sec:limits} concluded that the search for prior art should come first. Tested on the next
proposal, it did. We intended to apply $0$-dimensional sublevel persistent homology to the measured
dispersion: the roton as a homology class born at $\Delta_R$ and dying at the maxon $\Delta_M$, with
small-persistence pairs as instrumental noise. For a function on an interval this is exactly
topographic prominence --- as the founders of the subject state themselves --- and on the published
table the numbers agree to the digit: prominence $0.4501$~meV against $\Delta_M - \Delta_R =
0.4501$~meV. The persistence of the roton is a difference of two numbers already tabulated at every
pressure. Sublevel persistence of a dispersion with a noise threshold is in any case already
published for photonic band structures, and persistent homology of a Gross--Pitaevskii gas with
quantized vortices is published too --- which our own refuted vortex work of \S\ref{sec:empirical}
should have cited. Nothing was built. The episode cost a literature search rather than a release.

\section{Using and extending the library}

\paragraph{Reuse.} The modules marked in \S\ref{sec:library} are self-contained.
\lean{VortexWinding} and \lean{QuantizedCirculation} depend only on Mathlib and apply to any
code using plaquette phase winding. \lean{GPGalerkin} is stated over an arbitrary additive
group and finite mode set.

\paragraph{Open targets,} in the order we consider most valuable:
\begin{enumerate}[leftmargin=*,itemsep=2pt]
\item the sampling theorem linking the discrete loop sum to the continuum circulation for a
  $C^1$ phase with step bound below $\pi$---this would turn resolution criteria for vortex
  detection from folklore into theorem;
\item the oriented sum of plaquette windings over the faces of a grid cube vanishes (away
  from the $\pi$ case)---the discrete statement that vortex lines do not end, which is what
  threshold-free line tracing silently assumes;
\item the chain-rule step for \lean{GPGalerkin}, upgrading algebraic rates to conservation
  along solutions of the truncated system;
\item the Gross--Pitaevskii residual in the vocabulary of~\cite{OpenAINSE}, parallel to their
  Navier--Stokes residual, with the Madelung correspondence as a theorem between them;
\item the remaining coefficients $C,D,E,K,L$ of~\cite{Godfrin2021} Eq.~(22). The Bose integral is now
  proved at every natural order, so what remains is the Lagrange inversion $k(\omega)$ and the density
  of states --- combinatorial rather than analytic work. One structural point emerged in doing this:
  $D$ and $K$ carry $\zeta(7)$ and $\zeta(9)$, which have no known closed form, so those two
  coefficients can never be settled by exact algebra, only numerically. That is worth knowing before
  anyone tries to adjudicate the disagreement with earlier published series.
\end{enumerate}

\section{Conclusion}

We set out to find new structure in quantum fluids with a proof assistant and found, after
checking, that the structure we had was known. What remains is a library whose statements
are exactly as strong as they say, a small correctness result for a ubiquitous algorithm,
and a record in which every error is attributed to the mechanism that caught it. The record
suggests a division of labour: the kernel for truth, controls for measurement, a second
person for the build, and the library for novelty. We offer the first of these, with the
other three as its documented context.

\paragraph{Availability.} Source, claim ledger with all retractions
(\texttt{RETRACTIONS.md}), pre-registrations and run outputs:
\url{https://github.com/xaviercallens/SocrateAI-Scientific-QuantumFluids}.
The earlier report on the shell-model measurement programme remains in the repository at
tag \texttt{v1.1.0}.

\paragraph{Acknowledgements.} An external reviewer working from the LeanMaster project
rebuilt the tagged release independently and reported the two defects of
\S\ref{sec:caught}, items~5 and~7. The $^4$He dispersion data are those of Godfrin et
al.~\cite{Godfrin2021}; the turbulence snapshot is that of Polanco, M\"uller and
Krstulovic~\cite{PolancoData}.

\bibliographystyle{unsrt}
\bibliography{refs}
\end{document}
